% ICCSDI 2026 conference manuscript template
%
% International Conference on Computational Science and Data Intelligence
% (ICCSDI 2026), 11--12 December 2026, NMIMS Mumbai, India (Hybrid Mode)
%
% IMPORTANT
% 1. This package is intended only for ICCSDI 2026 conference-paper submissions.
% 2. Do not edit the supplied class or bibliography-style files.
% 3. Replace all instructional text and placeholders before submission.
% 4. Follow the current ICCSDI 2026 submission guidelines and page limits.

\documentclass[pdflatex,sn-mathphys-num]{sn-jnl}

%%%% Packages used by the ICCSDI 2026 conference template
\usepackage{graphicx}
\usepackage{multirow}
\usepackage{amsmath,amssymb,amsfonts}
\usepackage{amsthm}
\usepackage{mathrsfs}
\usepackage[title]{appendix}
\usepackage{xcolor}
\usepackage{textcomp}
\usepackage{manyfoot}
\usepackage{booktabs}
\usepackage{algorithm}
\usepackage{algorithmicx}
\usepackage{algpseudocode}
\usepackage{listings}

%%%% Optional theorem environments
\theoremstyle{thmstyleone}
\newtheorem{theorem}{Theorem}
\newtheorem{proposition}[theorem]{Proposition}

\theoremstyle{thmstyletwo}
\newtheorem{example}{Example}
\newtheorem{remark}{Remark}

\theoremstyle{thmstylethree}
\newtheorem{definition}{Definition}

\raggedbottom

\begin{document}

%%%% Manuscript title
% Provide a concise short title in square brackets.
\title[Short Manuscript Title]{Full Title of the Manuscript}

%%%% Authors and affiliations
% Mark the corresponding author with an asterisk.
% Add or remove author and affiliation blocks as required.
\author*[1]{\fnm{First} \sur{Author}}\email{corresponding.author@example.edu}

\author[2]{\fnm{Second} \sur{Author}}\email{second.author@example.edu}

\author[1,3]{\fnm{Third} \sur{Author}}\email{third.author@example.edu}

\affil*[1]{\orgdiv{Department/School}, \orgname{Institution Name},
  \orgaddress{\street{Street Address}, \city{City}, \postcode{Postal Code},
  \state{State}, \country{Country}}}

\affil[2]{\orgdiv{Department/School}, \orgname{Institution Name},
  \orgaddress{\street{Street Address}, \city{City}, \postcode{Postal Code},
  \state{State}, \country{Country}}}

\affil[3]{\orgdiv{Department/School}, \orgname{Institution Name},
  \orgaddress{\street{Street Address}, \city{City}, \postcode{Postal Code},
  \state{State}, \country{Country}}}

%%%% Abstract
% Write a self-contained abstract. State the problem, method, principal results,
% and implications. Avoid citations, undefined abbreviations, and equations.
\abstract{Provide a concise and self-contained abstract describing the research
problem, motivation, proposed approach, experimental or analytical methodology,
principal findings, and significance of the work. Clearly state the novelty
and original contribution of the manuscript. Replace this instructional text before submission.}

%%%% Keywords
% Use specific terms that support indexing and discovery.
\keywords{computational science, data intelligence, trustworthy artificial intelligence,
secure computational systems, keyword five}

\maketitle

\begin{center}
\small
\textbf{International Conference on Computational Science and Data Intelligence (ICCSDI 2026)}\\
11--12 December 2026, NMIMS Mumbai, India $\vert$ Hybrid Mode
\end{center}
\vspace{0.5em}

% The supplied class uses internal tables in the title block. Giving manuscript
% tables section-aware PDF anchors prevents duplicate hyperref destinations.
\renewcommand{\theHtable}{\thesection.\arabic{table}}

\section{Introduction}\label{sec:introduction}

Explain the research context, practical or theoretical motivation, and the
specific problem addressed. Identify the limitations of existing approaches
and state the research gap. Conclude the introduction with the objectives,
main contributions, and organisation of the manuscript.

A concise contribution list may be included when it improves clarity:
\begin{itemize}
\item State the first original technical or scientific contribution.
\item State the second contribution, including the evidence used to validate it.
\item State the third contribution and its significance or practical implication.
\end{itemize}

The manuscript should clearly establish its originality through appropriate
theoretical development, experiments, datasets, analyses, comparisons, or
applications. Cite related literature using the numbered citation
style, for example \cite{goodfellow2016deep,vaswani2017attention}.

\section{Related Work and Research Gap}\label{sec:related}

Critically review the most relevant and recent work. Organise the discussion by
methods, assumptions, application domains, datasets, or evaluation protocols
rather than listing papers sequentially. Explain how the present study differs
from and advances the state of the art.

\section{Materials and Methods}\label{sec:methods}

Describe the proposed method with enough detail to support verification and
reproducibility. Depending on the paper, this section may be renamed
``Proposed Methodology,'' ``System Model,'' ``Mathematical Framework,'' or
``Research Method.'' Define every symbol when it first appears.

\subsection{Problem Formulation}\label{subsec:problem}

State the inputs, outputs, assumptions, constraints, and optimisation or decision
objective. A general supervised learning objective can be written as
\begin{equation}
\theta^{\star}=\arg\min_{\theta}\left[
\frac{1}{n}\sum_{i=1}^{n}\mathcal{L}\bigl(f_{\theta}(\mathbf{x}_i),y_i\bigr)
+\lambda\,\Omega(\theta)\right],
\label{eq:objective}
\end{equation}
where $n$ is the number of observations, $\mathbf{x}_i$ and $y_i$ are the input
and target for observation $i$, $f_{\theta}$ is the model parameterised by
$\theta$, $\mathcal{L}$ is the task loss, $\Omega$ is a regularisation term, and
$\lambda\geq 0$ controls the regularisation strength. Replace this example with
the formulation used in the manuscript.

\subsection{Proposed Architecture or Workflow}\label{subsec:architecture}

Describe all components, data flows, interfaces, trust boundaries, and design
choices. Refer to each figure in the text before it appears. Place figure files
in the same folder as the manuscript and use PDF, PNG, or JPG with PDFLaTeX.

% Example figure block (uncomment and replace the file name):
% \begin{figure}[htbp]
% \centering
% \includegraphics[width=0.90\textwidth]{architecture.pdf}
% \caption{Overview of the proposed architecture. Explain all symbols and stages
% in the caption or the main text.}\label{fig:architecture}
% \end{figure}

\subsection{Algorithm or Computational Procedure}\label{subsec:algorithm}

Provide pseudocode when it helps readers reproduce the method. Define each
input and output, and discuss computational complexity where relevant.

\begin{algorithm}[htbp]
\caption{Generic training or optimisation procedure}\label{alg:procedure}
\begin{algorithmic}[1]
\Require Dataset $\mathcal{D}$, initial parameters $\theta_0$, maximum iterations $T$
\Ensure Trained parameters $\theta^{\star}$
\State $\theta \gets \theta_0$
\For{$t=1$ to $T$}
  \State Compute the objective or loss on the selected data
  \State Update $\theta$ using the stated optimisation rule
  \If{the stopping criterion is satisfied}
    \State \textbf{break}
  \EndIf
\EndFor
\State \Return $\theta^{\star}\gets\theta$
\end{algorithmic}
\end{algorithm}

\section{Experimental Design and Evaluation}\label{sec:experiments}

Provide enough information for independent reproduction. Report the dataset or
data source, inclusion and exclusion rules, preprocessing, train--validation--test
strategy, baselines, parameter settings, hardware and software environment,
random seeds, and statistical analysis. For security studies, clearly describe
the threat model, attacker capabilities, assets, assumptions, and evaluation
boundaries.

\subsection{Datasets and Preprocessing}\label{subsec:data}

Describe each dataset, its provenance, sample size, variables, licensing or
access conditions, preprocessing, missing-data treatment, and possible bias.
Do not disclose sensitive or personally identifiable information.

\subsection{Baselines and Evaluation Metrics}\label{subsec:metrics}

Justify the selected baselines and metrics. Report uncertainty using confidence
intervals, repeated runs, standard deviation, or suitable statistical tests when
appropriate. Avoid relying on a single metric when the application involves
class imbalance, safety, fairness, robustness, privacy, or security.

\subsection{Implementation Details}\label{subsec:implementation}

Report software libraries and versions, hardware, hyperparameters, stopping
criteria, and resource requirements. State whether code, models, prompts,
configuration files, and evaluation scripts will be made available.

\section{Results}\label{sec:results}

Present findings objectively and in the same order as the research questions or
hypotheses. Do not duplicate all table values in the prose. Use consistent
precision and clearly indicate whether higher or lower values are preferable.

\begin{table}[htbp]
\caption{Example comparison of the proposed method with representative baselines.}
\label{tab:results}
\centering
\begin{tabular}{@{}lccc@{}}
\toprule
Method & Metric 1 $\uparrow$ & Metric 2 $\uparrow$ & Cost $\downarrow$ \\
\midrule
Baseline A       & 0.812 & 0.774 & 1.00 \\
Baseline B       & 0.836 & 0.791 & 1.18 \\
Proposed method  & \textbf{0.861} & \textbf{0.824} & 1.12 \\
\botrule
\end{tabular}
\end{table}

Replace Table~\ref{tab:results} with verified results. Define all metrics,
units, symbols, arrows, boldface conventions, and statistical markers.

\section{Discussion}\label{sec:discussion}

Interpret the results in relation to the research questions and prior work.
Explain why the method succeeds or fails, identify practical implications, and
discuss generalisability. Include ablation, sensitivity, error, fairness,
robustness, privacy, security, or scalability analyses where relevant.

\subsection{Limitations and Threats to Validity}\label{subsec:limitations}

State methodological, data, measurement, implementation, and external-validity
limitations explicitly. Avoid claims that extend beyond the evidence.

\section{Conclusion and Future Work}\label{sec:conclusion}

Summarise the problem, approach, most important validated findings, and the
scientific or practical contribution. Do not introduce new results. End with
specific and realistic directions for future work.

\backmatter

\bmhead{Supplementary Information}

State whether supplementary files accompany the manuscript. Delete this heading
when it is not applicable or follow the conference submission instructions.

\bmhead{Acknowledgements}

Acknowledge non-author contributions, institutional support, facilities, and
other assistance. Keep this section brief and obtain permission before naming
individuals.

\section*{Declarations}

\textbf{Funding}\quad
State all funding sources and grant numbers, or write ``The authors received no
specific funding for this work.''

\textbf{Competing interests}\quad
State any financial or non-financial competing interests, or write ``The authors
declare no competing interests.''

\textbf{Ethics approval and consent to participate}\quad
Provide the approving body, approval number, and consent statement when human or
animal participants, samples, or identifiable data are involved. Otherwise write
``Not applicable.''

\textbf{Consent for publication}\quad
Provide the required consent statement, or write ``Not applicable.''

\textbf{Data availability}\quad
State where the data can be accessed, including a repository and persistent
identifier where possible. If data cannot be shared, explain the reason and the
conditions for controlled access.

\textbf{Materials availability}\quad
State the availability of specialised materials, or write ``Not applicable.''

\textbf{Code availability}\quad
State where the source code, trained models, prompts, configuration files, and
analysis scripts can be accessed, including version or release information.

\textbf{Author contributions}\quad
Describe each author's contribution using clear role-based statements. Ensure
that all listed authors satisfy recognised authorship criteria and the conference guidelines.

%%%% Optional appendix
% \begin{appendices}
% \section{Additional Derivations or Results}\label{app:additional}
% Place material that supports, but is not essential to, the main argument here.
% \end{appendices}

\bibliography{ICCSDI2026-references}

\end{document}
